% Contenidos del capítulo.
% Las secciones presentadas son orientativas y no representan
% necesariamente la organización que debe tener este capítulo.

\section{Introducción}


En el contexto actual, el desarrollo de tecnologías de realidad virtual (VR) ha generado una creciente demanda de herramientas y recursos que permitan a los desarrolladores crear experiencias inmersivas y funcionales. Uno de los desafíos más significativos dentro de este ámbito es proporcionar soluciones prácticas y accesibles para la interacción en entornos virtuales, especialmente en lo que respecta a la entrada de texto.

Los teclados virtuales son un componente fundamental para muchas aplicaciones de VR, ya que permiten a los usuarios realizar tareas como introducir texto, realizar búsquedas, o interactuar con interfaces complejas sin necesidad de hardware adicional. Sin embargo, las opciones disponibles en el mercado suelen estar limitadas en términos de funcionalidades, idiomas soportados y accesibilidad, lo que deja un vacío importante que podría ser abordado mediante un enfoque más inclusivo y personalizable.

\section{Motivación}
La idea de este Asset me surgió durante la realización de las prácticas curriculares en el Instituto de Robótica y Tecnologías de la Información y la Comunicación(IRTIC) donde se me pidió un teclado virtual para uno de los proyectos que se estaban realizando con Unity.

Al buscar en el Asset Store (tienda oficial de Unity), no encontré nada que se adaptara tanto a las necesidades idiomáticas como presupuestarias. Por lo que el equipo decidió hacer un teclado a mano muy básico.

Tras la realización de este, uno de los tutores de la empresa, que también resultaba ser profesor de la universidad, me sugirió la idea de hacerlo mi TFG, puesto que, para entonces aun no tenía una idea clara de qué elegir como producto de mi TFG. Pero añadirle mas funcionalidades, como el diseñarlo para varios idiomas para captar otros mercados que también sufren el hecho de tener que caer en la lengua anglosajona.

\section{Objetivos}

Este trabajo de fin de grado se centra en el diseño y desarrollo del Asset de un teclado multilingüe para entornos de realidad virtual utilizando la plataforma Unity. La propuesta, no solo busca satisfacer las necesidades específicas detectadas durante las prácticas en el Instituto de Robótica y Tecnologías de la Información y la Comunicación (IRTIC), sino también ofrecer una herramienta competitiva y versátil que pueda ser utilizada por desarrolladores y usuarios en una amplia variedad de contextos.

Entre las funcionalidades más destacadas se incluyen el soporte para múltiples idiomas, con la finalidad de que el usuario no tenga que hacer mucho a la hora de instaurar el que desee.


\section{Organización de la memoria}

La memoria se organizará de la siguiente manera:

\begin{itemize}
	\item \textbf{Estado del Arte (Capítulo 2):} en este punto realiza una revisión exhaustiva de las tecnologías utilizadas para la realización del teclado y de las aplicaciones parecidas.
	\item \textbf{Requisitos, especificaciones, coste, riesgos, viabilidad (Capítulo 3):} en este apartado se abordan los aspectos clave del proyecto relacionados con la planificación,  la fundamentación, detalles técnicos de la implementación, los costes para llevar a cabo éste, identificación y planes de acción de riesgos que puedan surgir y un análisis técnico, económico y operativo para comprender si se puede llevar a cabo el proyecto de una manera sostenible. 
	\item \textbf{Análisis (Capítulo 4):} en esta parte del TFG se realiza una revisión exhaustiva de las tecnologías utilizadas para la realización del teclado. 
	\item \textbf{Diseño (Capítulo 5):} la realización de una descripción detallada de los distintos diseños utilizados como la intefaz de usuarios, algoritmos utilizados y funcionalidades específicas es el centro de estudio de este capítulo.
	\item \textbf{Implementación y pruebas (Capítulo 6):} este punto del trabajo comporta la realización de una revisión exhaustiva de las tecnologías utilizadas para la realización del teclado y de las distintas pruebas que se llevarán a cabo.
	\item \textbf{Conclusiones (Capítulo 7):} se resumen los logros alcanzados, las implicaciones prácticas y teóricas del teclado diseñado, y se proponen direcciones para futuras actualizaciones.
	\item \textbf{Apéndice :} Aquí se adjunta material adicional, como el código fuente, datos adicionales y gráficos detallados que respaldan el contenido principal de la memoria. 
	\item \textbf{Bibliografía:} incluye una lista completa de todas las fuentes citadas en el trabajo.
\end{itemize}
